\documentclass[10pt,letterpaper,oneside]{report}
\usepackage[T1]{fontenc}
\usepackage[utf8]{inputenc}
\usepackage{lmodern}
\usepackage{microtype}
\usepackage[margin=0.85in,headheight=15pt]{geometry}
\usepackage{enumitem,booktabs,longtable,tabularx,array,makecell}
\usepackage{fancyhdr,lastpage,xcolor,hyperref,titlesec}
\usepackage{listings}
\usepackage[most]{tcolorbox}
\definecolor{ink}{HTML}{182D40}
\definecolor{accent}{HTML}{126B77}
\definecolor{pale}{HTML}{EEF5F5}
\definecolor{grayblue}{HTML}{677887}
\definecolor{codebg}{HTML}{F5F7F8}
\hypersetup{colorlinks=true,linkcolor=accent,urlcolor=accent,pdfauthor={C23 Engineering Companion},pdfsubject={ISO/IEC 9899:2024 educational engineering companion}}
\titleformat{\chapter}[display]{\normalfont\huge\bfseries\color{ink}}{\filleft\large\color{accent}CHAPTER \thechapter}{1ex}{\titlerule[1pt]\vspace{1ex}}
\titleformat{\section}{\Large\bfseries\color{ink}}{\thesection}{0.7em}{}
\titleformat{\subsection}{\large\bfseries\color{accent}}{\thesubsection}{0.65em}{}
\setlist[itemize]{leftmargin=1.3em,itemsep=2pt,topsep=3pt}
\setlist[enumerate]{leftmargin=1.6em,itemsep=3pt,topsep=3pt}
\setlength{\parskip}{5pt}
\setlength{\parindent}{0pt}
\renewcommand{\arraystretch}{1.18}
\pagestyle{fancy}\fancyhf{}\fancyhead[L]{\small\sffamily C23 ENGINEERING SERIES}\fancyhead[R]{\small\sffamily\color{accent}\DocShort}\fancyfoot[L]{\small\color{grayblue}ISO/IEC 9899:2024 companion}\fancyfoot[R]{\small\thepage\ / \pageref{LastPage}}
\newtcolorbox{keybox}{breakable,colback=pale,colframe=accent,boxrule=.6pt,arc=1mm,left=8pt,right=8pt,top=6pt,bottom=6pt}
\newtcolorbox{warnbox}{breakable,colback=codebg,colframe=grayblue,boxrule=.6pt,arc=1mm,left=8pt,right=8pt,top=6pt,bottom=6pt}
\lstdefinestyle{cstyle}{language=C,basicstyle=\ttfamily\footnotesize,keywordstyle=\color{accent}\bfseries,commentstyle=\color{grayblue}\itshape,stringstyle=\color{ink},backgroundcolor=\color{codebg},frame=single,rulecolor=\color{grayblue!40},breaklines=true,showstringspaces=false,tabsize=4,columns=fullflexible,keepspaces=true,upquote=true}
\lstset{style=cstyle}
\newcommand{\refcl}[1]{\textsf{\footnotesize[ISO/IEC 9899:2024, #1]}}
\newcommand{\standard}{\textbf{[STANDARD]}}
\newcommand{\optional}{\textbf{[OPTIONAL]}}
\newcommand{\guidance}{\textbf{[GUIDANCE]}}
\newcommand{\testtag}{\textbf{[TEST]}}
\newcommand{\term}[1]{\textbf{#1}}
\newcommand{\DocShort}{LIBRARY REFERENCE}
\begin{document}
\begin{titlepage}
\color{ink}\vspace*{1.2in}
{\sffamily\small\color{accent}\bfseries TECHNICAL DOCUMENTATION SUITE \hfill EDITION 2026-10}\par
\vspace{.3in}\rule{\linewidth}{1.5pt}\par\vspace{.22in}
{\Huge\bfseries C23 Standard Library Reference}\par\vspace{.22in}
{\Large\color{grayblue}Header catalog, contracts, failure modes and C23 library additions}\par\vspace{.4in}
\begin{keybox}\textbf{Standards basis:} ISO/IEC 9899:2024 (C23). Independent educational and engineering companion. Not an official ISO interpretation, full specification, certification or conformance assessment.\end{keybox}
\vfill{\large\bfseries C23 Software Engineering and Language Standards}\par\vspace{.15in}{\color{grayblue}Prepared October 2026 \quad | \quad Source edition: October 2024}\par
\end{titlepage}
\pagenumbering{roman}\tableofcontents\clearpage\pagenumbering{arabic}

\chapter*{Sources and interpretation key}
\addcontentsline{toc}{chapter}{Sources and interpretation key}
\textbf{Primary source:} ISO/IEC 9899:2024(en), \emph{Information technology --- Programming languages --- C}, fifth edition (2024). Clause and annex identifiers refer to the supplied licensed standard. This handbook contains original explanations, designs and examples; it does not reproduce the standard. The licensed standard remains the authoritative source for exact wording and conformance decisions.\par
\textbf{Secondary engineering sources:} \url{https://gcc.gnu.org/onlinedocs/gcc/Standards.html}; \url{https://clang.llvm.org/c_status.html}; \url{https://cmake.org/cmake/help/latest/prop_tgt/C_STANDARD.html}. Compiler support and diagnostics are release-specific: verify the installed compiler, target and library before relying on a feature.\par
\begin{keybox}\textbf{Interpretation categories.} \textbf{[STANDARD]} summarizes a requirement, constraint, semantic rule or defined facility and cites the applicable clause. \textbf{[OPTIONAL]} describes conditional or implementation-dependent support. \textbf{[GUIDANCE]} records an engineering recommendation, not an ISO mandate. \textbf{[TEST]} is an original check or exercise; a passing test is evidence but is not proof of full conformity.\end{keybox}
\textbf{Copyright boundary:} The supplied PDF is individually licensed. Do not redistribute it with this suite; distribute only these independently written companion materials in accordance with applicable rights.\par


\chapter{Library availability model}
The standard library is not a platform-independent guarantee that every target exports every interface. Hosted versus freestanding rules and conditional features matter. The library clause is normative for applicable facilities and implementations; test APIs at compile, link and runtime stages when necessary. \refcl{5.2.2, 7.1}.
\section{Headers at a glance}
\scriptsize\begin{longtable}{p{.18\linewidth}p{.075\linewidth}p{.26\linewidth}p{.375\linewidth}}\toprule Header & Clause & Purpose & Audit focus\\\midrule\endhead
\texttt{<assert.h>} & 7.2 & Assertions and diagnostics & Do not use assertions as the sole validation of untrusted input.\\\hline
\texttt{<complex.h>} & 7.3 & Complex arithmetic & Review branch cuts and floating assumptions.\\\hline
\texttt{<ctype.h>} & 7.4 & Byte character classification & Convert possibly negative char to unsigned char or EOF before calls.\\\hline
\texttt{<errno.h>} & 7.5 & Error reporting & Inspect immediately when the associated API specifies it.\\\hline
\texttt{<fenv.h>} & 7.6 & Floating-point environment & Observe floating environment access requirements.\\\hline
\texttt{<float.h>} & 7.7 & Floating type characteristics & Do not assume IEEE formats without evidence.\\\hline
\texttt{<inttypes.h>} & 7.8 & Integer format macros & Keep format types aligned with arguments.\\\hline
\texttt{<iso646.h>} & 7.9 & Alternative operator spellings & Compatibility header; avoid unnecessary dependencies.\\\hline
\texttt{<limits.h>} & 7.10 & Integer type limits & Size algorithms from limits rather than fixed guesses.\\\hline
\texttt{<locale.h>} & 7.11 & Localization & Avoid uncontrolled process-global locale changes.\\\hline
\texttt{<math.h>} & 7.12 & Mathematics & Review domain/range errors and floating exceptions.\\\hline
\texttt{<setjmp.h>} & 7.13 & Nonlocal jumps & Avoid crossing invalid lifetimes or threads.\\\hline
\texttt{<signal.h>} & 7.14 & Signal handling & Use only permitted operations in signal contexts.\\\hline
\texttt{<stdalign.h>} & 7.15 & Alignment compatibility & Prefer unambiguous C23 spellings where compatible.\\\hline
\texttt{<stdarg.h>} & 7.16 & Variadic arguments & Ensure default promotions and format consistency.\\\hline
\texttt{<stdatomic.h>} & 7.17 & Atomic types and operations & Use coherent memory-order reasoning.\\\hline
\texttt{<stdbit.h>} & 7.18 & C23 bit and byte operations & Compile-probe type-generic macros and implementation support.\\\hline
\texttt{<stdbool.h>} & 7.19 & Boolean compatibility & Avoid collisions with newly standardized keywords.\\\hline
\texttt{<stdckdint.h>} & 7.20 & C23 checked arithmetic & Check overflow return and output semantics.\\\hline
\texttt{<stddef.h>} & 7.21 & Core definitions and nullptr\_t & Distinguish nullptr from integer zero.\\\hline
\texttt{<stdint.h>} & 7.22 & Integer-width types & Availability of exact widths follows defined conditions.\\\hline
\texttt{<stdio.h>} & 7.23 & Streams and formatted I/O & Check return values, widths and format correspondence.\\\hline
\texttt{<stdlib.h>} & 7.24 & Allocation, conversion and utilities & Preserve ownership and handle allocation failure.\\\hline
\texttt{<stdnoreturn.h>} & 7.25 & Legacy no-return macro & Review migration to noreturn attribute.\\\hline
\texttt{<string.h>} & 7.26 & Byte strings and memory & Differentiate byte count, object extent and terminator.\\\hline
\texttt{<tgmath.h>} & 7.27 & Type-generic math & Check chosen overload and floating types.\\\hline
\texttt{<threads.h>} & 7.28 & Threading primitives & Plan mutex, lifetime and shutdown invariants.\\\hline
\texttt{<time.h>} & 7.29 & Time and dates & Check timescale, conversion and range.\\\hline
\texttt{<uchar.h>} & 7.30 & Unicode conversions & Track conversion state and encoding errors.\\\hline
\texttt{<wchar.h>} & 7.31 & Wide character utilities & Handle units and conversion state.\\\hline
\texttt{<wctype.h>} & 7.32 & Wide character classification & Use locale-conscious classification.\\\hline
\bottomrule\end{longtable}\normalsize
\chapter{Contracts and common failure modes}
\section{Standard library contract and formatted I/O}
\standard\ For each call distinguish preconditions, diagnostics, effects, error returns, and stream orientation. Treat return counts as control data, not decorations. \refcl{7.1 and 7.23}.\par\guidance\ Use snprintf with a known destination size; inspect truncation; do not use untrusted data as the format string.

\section{ctype input domain}
\standard\ Character classification takes values representable as unsigned char or EOF. A signed char with a negative value is not safe to pass directly. \refcl{7.4}.\par\guidance\ Use isalpha((unsigned char)c) only when c is an actual char-valued byte; preserve EOF as a separate case.

\section{Floating-point algorithms}
\standard\ The observable behavior of floating exceptions, rounding and math errors depends on applicable language and annex provisions. \refcl{7.6, 7.7, 7.12}.\par\guidance\ Document NaN, infinity, signed zero and permitted rounding when comparing numeric results.

\section{Atomics and thread interfaces}
\standard\ C atomics provide specified operations and memory orders; C threads provide mutexes, condition variables and thread lifecycle facilities. \refcl{7.17, 7.28}.\par\guidance\ Prefer mutex-protected invariants first; annotate happens-before claims with release/acquire or lock edges.

\section{std(bit) integer bit utilities}
\standard\ Type-generic bit operations include counts of leading or trailing bits and width/floor/ceil routines. \refcl{7.18}.\par\guidance\ Do not substitute floating logarithms for exact integer bit-width operations.

\section{stdckdint checked arithmetic}
\standard\ The ckd\_add, ckd\_sub and ckd\_mul facilities detect overflow of the result type and provide an output through a pointer. \refcl{7.20}.\par\guidance\ Check the indicated overflow; do not subsequently consume overflowed outputs without a documented policy.

\section{nullptr\_t and unreachable}
\standard\ C23 offers a null pointer type and an unreachable optimization hint; reaching unreachable is undefined behavior. \refcl{7.21}.\par\guidance\ Never use unreachable to hide a real unhandled error branch.

\section{scanf and printf formats}
\standard\ Conversion specifications must correspond to the supplied argument types; widths and buffer sizes matter. \refcl{7.23}.\par\guidance\ Prefer fgets plus explicit parsing over unconstrained string scans.

\section{Allocation and lifetime}
\standard\ Allocation routines return storage subject to alignment, ownership and failure; frees end lifetime. realloc with zero size is undefined in C23. \refcl{7.24}.\par\guidance\ Use a temporary pointer for realloc; define zero-sized allocation handling and secure erasure separately.

\section{memcpy, memmove, memset\_explicit}
\standard\ memcpy requires nonoverlapping regions; memmove is for possibly overlapping regions. memset\_explicit supports explicit wiping. \refcl{7.26}.\par\guidance\ Prove both buffer extents; treat secret erasure as only one part of secret lifetime control.

\section{timegm, timespec\_getres and time conversions}
\standard\ Time bases, UTC conversion and range behavior require testing with the implementation. \refcl{7.29}.\par\guidance\ Define time-zone policy and verify leap/overflow boundary tests.

\section{UTF and wide-character facilities}
\standard\ Conversions depend on input validity and often on conversion state. \refcl{7.30--7.32}.\par\guidance\ Keep encoding and unit counts explicit; invalid input should return an error rather than silently truncate.

\chapter{Original C23 library implementation examples}
\section{Checked multiplication for array allocation}
\begin{lstlisting}
#include <stddef.h>
#include <stdlib.h>
#include <stdckdint.h>
void *allocate_array(size_t count, size_t element_size) {
    size_t total;
    if (ckd_mul(&total, count, element_size)) return NULL;
    if (total == 0) return NULL; /* project policy */
    return malloc(total);
}
\end{lstlisting}
\guidance\ Checked arithmetic can prevent a wrapped allocation size.\par\refcl{7.20, 7.24.4}.\optional\ Compile only on implementations that supply \texttt{<stdckdint.h>} and the invoked macros. For a portable fallback, check \texttt{count > SIZE\_MAX / element\_size} after verifying the divisor is nonzero.
\section{Byte classification without signed-char misuse}
\begin{lstlisting}
#include <ctype.h>
#include <stdbool.h>
bool ascii_decimal_byte(char value) {
    unsigned char u = (unsigned char)value;
    return u >= (unsigned char)'0' && u <= (unsigned char)'9';
}
\end{lstlisting}
\guidance\ This example deliberately avoids locale-dependent decimal character classification; \texttt{isdigit} is appropriate where its specified domain and semantics meet the use case.
\section{Bounded string formatting}
\begin{lstlisting}
#include <stdio.h>
#include <stdbool.h>
bool make_label(char *dst, size_t size, unsigned id) {
    if (dst == NULL || size == 0) return false;
    int count = snprintf(dst, size, "item-%u", id);
    return count >= 0 && (size_t)count < size;
}
\end{lstlisting}
\section{Replacing an overlapping copy}
\begin{lstlisting}
#include <string.h>
void erase_first_byte(char *buf, size_t bytes) {
    if (bytes > 1) memmove(buf, buf + 1, bytes - 1);
}
\end{lstlisting}
\guidance\ The caller guarantees \texttt{buf} points to at least \texttt{bytes} bytes. Using \texttt{memcpy} for overlapping ranges violates its preconditions.
\chapter{Library verification worksheet}
For each library facility record: required header; clause reference; feature-test macro if any; hosted/freestanding applicability; compiler version; libc version; compile/link/run probe; preconditions; postconditions; error handling; thread/locale interactions; tests; fallback; owner.\par
\begin{keybox}\textbf{Minimum acceptance gate [GUIDANCE]:} a successful header probe, a semantic compile probe, a link probe for non-inline functions, a failure-path test and a documented target support decision.\end{keybox}
\chapter{Annex and library cross-reference}
\begin{itemize}
\item Annex B is an informative navigation summary of library facilities, not an alternate normative contract.
\item Annex F--H affect floating-point behavior when applicable, and their conditions must be read carefully.
\item Annex K specifies optional bounds-checking interfaces; it is normative for implementations claiming that support, not mandatory on all platforms.
\item Annex L likewise imposes requirements conditionally where the implementation declares analyzability support.
\item Annex M is an informative change index for fifth-edition library additions.
\end{itemize}

\end{document}
